\documentclass[10pt]{article}
\usepackage[a4paper,margin=0.3in]{geometry}
\usepackage{etoolbox}
\AtBeginEnvironment{align}{\small}
\usepackage{amsmath,amssymb,amsfonts,mathtools}
\usepackage[dvipsnames]{xcolor}
\usepackage{array,booktabs,float,longtable}
\usepackage[colorlinks=true,linkcolor=blue]{hyperref}
\allowdisplaybreaks
\numberwithin{equation}{section}
\setlength{\jot}{3pt}
\setlength{\emergencystretch}{3em}
\definecolor{JIM}{RGB}{0,130,60}
\definecolor{CC}{RGB}{140,0,170}
\definecolor{BFK}{RGB}{215,95,0}
\newcommand{\new}[1]{{\color{blue}#1}}
\newcommand{\old}[1]{{\color{red}#1}}
\newcommand{\nc}{\color{blue}}
\newcommand{\rc}{\color{red}}
\newcommand{\NN}{\mathcal N\log\frac{k^+}{\Lambda}}
\newcommand{\NNv}{\mathcal N\log\frac{\vee}{k^+}}
\newcommand{\Lv}{\log{\left(\frac{\vee}{k^{+}}\right)}}
\def\mf#1{\mathbb{#1}}

\newcommand{\Kc}{\mathcal K}
\title{Step by step from your equations:\\ where the extra $\log\frac{\vee}{k^+}$ terms of Group III Rows III, V, VI and Group II Row IV come from;\\ and the $\mu$ bookkeeping (running coupling, DGLAP)}
\date{}
\begin{document}
\maketitle
\noindent Your notation throughout. Your plus prescription:
\begin{equation}
\int dp^+\left[\frac{1}{p^+}\right]_+f(p^+)\equiv\int dp^+\,\frac{f(p^+)-f(0)}{p^+},\qquad
\int\frac{d\xi}{\xi^2}\,\frac{\xi}{(1-\xi)_+}\,g(\xi)\equiv\int\frac{d\xi}{\xi(1-\xi)}\,\big[g(\xi)-g(1)\big].
\label{eq:plus}
\end{equation}
Your \texttt{BFKL.tex} calls the Group III rows ``III, IV, V''; in \texttt{Group\_III.tex} they are Rows III, V, VI (Row IV has only four $\rho$'s). Equation numbers ``your (1.x)'' are those of your files as they compile.

\paragraph{Your ``first part of the calculation''.} The \texttt{.tex} file you mention did not reach me: the last files I received are \texttt{Group\_I/II/III.tex} and \texttt{4rho.tex} (2.30.24 AM). Please upload it again and I will check it against the three Groups.

\section*{Part A. The extra terms in line 1}
\section{The answer in one paragraph}
Nothing is added to your cross section. Every ``extra'' term is a piece of your line 1, the $p^+$ (resp.\ $\xi$) integral that runs up to $\vee-k^+$ (resp.\ down to $\frac{k^+}{\vee}$). It is written out because it grows like $\log\frac{\vee}{k^+}$. There are three of them in each row:
\begin{description}
\item[(i)] The $\log\frac{\vee}{k^+}$ in your line 2, from $\log(\frac{\vee}{k^+}-1)=\log\frac{\vee}{k^+}+\log(1-\frac{k^+}{\vee})$.
\item[(ii)] Minus the same, hidden in the subtraction $-f(0)$ of $\left[\frac{1}{p^+}\right]_+$ in line 1. (i) and (ii) cancel exactly. So the $\log\frac{\vee}{k^+}$ you see in line 2 is \emph{not} a $\vee$ dependence of the row: before your plus prescription, Eq.~(your 1.12), there is no such logarithm (Step 3).
\item[(iii)] The term (B), hidden in the $-\frac{1}{k^+}\delta_{im}\delta_{jk'}$ term of line 1. This term has no $\frac{1}{p^+}$, but after the $y'$ integral it becomes \emph{exactly} $\frac{k^+}{p^++k^+}\times$ (a $p^+$-independent integral) $+$ (a part that falls faster). Then $\int_0^{\vee-k^+}\frac{dp^+}{p^++k^+}=\log\frac{\vee}{k^+}$ (Step 4). This is the genuine BFKL term of the row.
\end{description}
In Group II Row IV, (i) and (ii) are the same with $\xi$. (iii) is $H_4$, from the $\log(y'-y)^2$ of the pole bracket in the $\frac{1-\xi}{\xi}$ term (Sec.~\ref{sec:G2IV}).

\section{Group III Row V, step by step}\label{sec:V}
\paragraph{Step 1: your equation after the $w'$ integral.}
\begin{align}
     \frac{d^{3}N_{\rm NLO}}{d^{3}k} & = \frac{1}{(2\pi)^{3}}\, \frac{ig^{4}\,f^{c'd'a}}{4\pi^{4}}\, \frac{1}{k^{+}}\,\int_{\Lambda}^{\vee-k^{+}}\, \frac{dp^{+}}{2\pi}\, \int_{w,y',z,y,x,x'}\, e^{-ik\cdot (y'-w)}\, e^{-ik\cdot \frac{p^{+}}{k^{+}}\,(y'-z)}\,\frac{(y'-z)^{m}\,(x'-y')^{k'}\,(y-w)^{i}\,(x-z)^{j}}{(y'-z)^{2}\,(x'-y')^{2}\,(y-w)^{2}\,(x-z)^{2}}\notag \\ & \left[\frac{\delta_{k'm}\,\delta_{ij}}{(p^{+}+k^{+})}-\frac{1}{p^{+}}\,\delta_{jm}\,\delta_{ik'}-\frac{1}{k^{+}}\, \delta_{im}\, \delta_{jk'}\right]\, \left[U^{e'c'}(y')\,U^{d'b}(z)\,\rho^{e'}(x')-U^{bd'}(z)\,U^{c'e'}(x')\,\rho^{e'}(x')\right]\times\notag \\ &  \left[U^{bc}(z)\,U^{ce}(x)\,U^{ad}(y)\,\rho^{e}(x)\rho^{d}(y)-U^{ad}(w)\,U^{ce}(x)\,U^{cb}(z)\,\rho^{e}(x)\rho^{d}(y)\right]\tag{your (1.12)}
\end{align}

After the Wilson lines are simplified (your Eq.~(1.14)), the colour factor is
\begin{equation}
\mathcal C(y')=\Big[U^{e'c'}(y')\,\rho^{e'}(x')-U^{c'e'}(x')\,\rho^{e'}(x')\Big]\times\Big[U^{ce}(x)\,U^{ad}(y)\,\rho^{e}(x)\rho^{d}(y)-U^{ad}(w)\,U^{ce}(x)\,\rho^{e}(x)\rho^{d}(y)\Big],
\label{eq:C}
\end{equation}
with $f^{c'ca}$ and $C\equiv\frac{1}{(2\pi)^{3}}\frac{ig^{4}f^{c'ca}}{4\pi^{4}}\frac{1}{k^{+}}$ in front. Only $U^{e'c'}(y')$ depends on $y'$.

\paragraph{Step 2: names.} Contract the Kronecker deltas of the three terms of the bracket. For each $p^+$,
\begin{align}
F_a(p^+)&=\int_{w,y',z,y,x,x'}e^{-ik\cdot (y'-w)}\, e^{-ik\cdot \frac{p^{+}}{k^{+}}(y'-z)}\,\frac{(y'-z)\cdot (x'-y')\,(y-w)\cdot (x-z)}{(y'-z)^{2}(x'-y')^{2}(y-w)^{2}(x-z)^{2}}\,\mathcal C(y') &&(\delta_{k'm}\delta_{ij}),\notag\\
F_b(p^+)&=\int_{w,y',z,y,x,x'}e^{-ik\cdot (y'-w)}\, e^{-ik\cdot \frac{p^{+}}{k^{+}}(y'-z)}\,\frac{(y'-z)\cdot (x-z)\,(x'-y')\cdot (y-w)}{(y'-z)^{2}(x'-y')^{2}(y-w)^{2}(x-z)^{2}}\,\mathcal C(y') &&(\delta_{jm}\delta_{ik'}),\notag\\
F_c(p^+)&=\int_{w,y',z,y,x,x'}e^{-ik\cdot (y'-w)}\, e^{-ik\cdot \frac{p^{+}}{k^{+}}(y'-z)}\,\frac{(y'-z)\cdot (y-w)\,(x'-y')\cdot (x-z)}{(y'-z)^{2}(x'-y')^{2}(y-w)^{2}(x-z)^{2}}\,\mathcal C(y') &&(\delta_{im}\delta_{jk'}).
\label{eq:F}
\end{align}
Your Eq.~(1.12), with the Wilson lines simplified, is then
\begin{equation}
\frac{d^3N_{\rm V}}{d^3k}=C\int_{\Lambda}^{\vee-k^{+}}\frac{dp^{+}}{2\pi}\bigg[\frac{F_a(p^+)}{p^++k^+}-\frac{F_b(p^+)}{p^+}-\frac{F_c(p^+)}{k^+}\bigg].
\label{eq:start}
\end{equation}

\paragraph{Step 3: your plus prescription, and what it does to $\vee$.} In your notation, exactly,
\begin{equation}
\int_\Lambda^{\vee-k^+}dp^+\,\frac{F_b(p^+)}{p^+}=\int_\Lambda^{\vee-k^+}dp^+\left[\frac{1}{p^+}\right]_+F_b(p^+)+F_b(0)\int_\Lambda^{\vee-k^+}\frac{dp^+}{p^+},\qquad
\int_\Lambda^{\vee-k^+}\frac{dp^+}{p^+}=\log\frac{k^+}{\Lambda}+\log\Big(\frac{\vee}{k^+}-1\Big).
\label{eq:step3}
\end{equation}
With $\int_\Lambda\to\int_0$ in the plus-prescribed term ($O(\Lambda)$), Eq.~\eqref{eq:start} is your Eq.~(1.14):
\begin{align}
\text{line 1}&=C\int_{0}^{\vee-k^{+}}\frac{dp^{+}}{2\pi}\bigg[\frac{F_a(p^+)}{p^++k^+}-\frac{F_b(p^+)-F_b(0)}{p^+}-\frac{F_c(p^+)}{k^+}\bigg],\notag\\
\text{line 2}&=-\frac{C}{2\pi}\,F_b(0)\,\Big[\log\frac{k^+}{\Lambda}+\log\Big(\frac{\vee}{k^+}-1\Big)\Big].
\label{eq:lines}
\end{align}
($\frac{C}{2\pi}=\frac{1}{(2\pi)^3}\frac{ig^4f^{c'ca}}{8\pi^5}\frac{1}{k^+}$, and $F_b(0)$, without the $p^+$ phase, is the integral of your line 2.)

\textbf{The point.} The left side of Eq.~\eqref{eq:step3} does \emph{not} grow like $\log\vee$: its integrand $\frac{F_b(p^+)}{p^+}$ falls at least like $\frac{1}{p^{+2}}$ (Step 4). The right side has $+F_b(0)\log\frac{\vee}{k^+}$ explicitly. So the plus-prescribed integral must contain $-F_b(0)\log\frac{\vee}{k^+}$: it subtracts $\frac{F_b(0)}{p^+}$ all the way up to $p^+=\vee-k^+$. Exactly: since $\frac{1}{p^+}=\frac{1}{p^++k^+}+\frac{k^+}{p^+(p^++k^+)}$ and $\int_0^{\vee-k^+}\frac{dp^+}{p^++k^+}=\log\frac{\vee}{k^+}$,
\begin{equation}
\int_0^{\vee-k^+}dp^+\,\frac{F_b(p^+)-F_b(0)}{p^+}=\int_0^{\vee-k^+}dp^+\bigg[\frac{k^+\big(F_b(p^+)-F_b(0)\big)}{p^+(p^++k^+)}+\frac{F_b(p^+)}{p^++k^+}\bigg]\;-\;F_b(0)\,\log\frac{\vee}{k^+}.
\label{eq:I2}
\end{equation}
This is identity (I2). Both integrands on the right fall like $\frac{1}{p^{+2}}$, so they are finite for $\vee\to\infty$. Put Eq.~\eqref{eq:I2} into line 1, Eq.~\eqref{eq:lines}:
\begin{align}
\text{(i)}\;&=-\frac{C}{2\pi}\,F_b(0)\log\frac{\vee}{k^+}\qquad\text{(the }\log\tfrac{\vee}{k^+}\text{ of line 2)},\notag\\
\text{(ii)}\;&=+\frac{C}{2\pi}\,F_b(0)\log\frac{\vee}{k^+}\qquad\text{(from }-\left[\tfrac{1}{p^+}\right]_+\text{ in line 1)},\qquad \text{(i)}+\text{(ii)}=0.
\label{eq:i_ii}
\end{align}
Written out, (i) is the first and (ii) the second $\log\frac{\vee}{k^+}$ line of Group III Row V in \texttt{Checks\_JIMWLK\_BFKL.pdf}.

\paragraph{Step 4: the $y'$ integral, exactly.} The two phases combine to one:
\begin{equation}
e^{-ik\cdot (y'-w)}\,e^{-ik\cdot\frac{p^+}{k^+}(y'-z)}=e^{-ik\cdot(z-w)}\,e^{-iq\cdot(y'-z)},\qquad q=t\,k,\quad t=\frac{p^++k^+}{k^+}.
\end{equation}
In each $F_n$ the integrand is $\frac{(y'-z)^m}{(y'-z)^2}G_n^m(y')$, where $G_n(y')$ collects $\frac{(x'-y')^{k'}}{(x'-y')^2}$ and $\mathcal C(y')$. Write, identically,
\begin{equation}
G_n(y')=G_n(z)+\big[G_n(y')-G_n(z)\big].
\label{eq:split}
\end{equation}
For the $G_n(z)$ part the $y'$ integral is a pure 2-D Fourier transform. With $\int d^2r\,e^{-iq\cdot r}\frac{r^m}{r^2}=-2\pi i\frac{q^m}{q^2}$ (from $\frac{r^m}{r^2}=\frac12\partial^m\log r^2$ and $\int d^2r\,e^{-iq\cdot r}\log r^2=-\frac{4\pi}{q^2}$):
\begin{equation}
\int_{y'}e^{-iq\cdot(y'-z)}\frac{(y'-z)^m}{(y'-z)^2}=-2\pi i\,\frac{1}{t}\,\frac{k^m}{k^2}=\frac{k^+}{p^++k^+}\int_{w'}e^{-ik\cdot(w'-z)}\frac{(w'-z)^m}{(w'-z)^2}.
\label{eq:I3}
\end{equation}
This is identity (I3). It is exact; no large-$p^+$ limit is taken. Hence, for $n=a,b,c$,
\begin{equation}
F_n(p^+)=\frac{k^+}{p^++k^+}\,F_n^{z}+F_n^{\Delta}(p^+),
\label{eq:Fsplit}
\end{equation}
with
\begin{align}
F_c^{z}&=\int_{w,w',z,y,x,x'}e^{-ik\cdot (w'-w)}\,\frac{(w'-z)\cdot (y-w)\,(x'-z)\cdot (x-z)}{(w'-z)^{2}(y-w)^{2}(x'-z)^{2}(x-z)^{2}}\,\mathcal C(z),\notag\\
F_c^{\Delta}(p^+)&=\int_{w,y',z,y,x,x'}e^{-ik\cdot (y'-w)}e^{-ik\cdot\frac{p^+}{k^+}(y'-z)}\frac{(y'-z)^m(y-w)^m}{(y'-z)^2(y-w)^2}\,\frac{(x-z)^{k'}}{(x-z)^2}\bigg[\frac{(x'-y')^{k'}}{(x'-y')^2}\mathcal C(y')-\frac{(x'-z)^{k'}}{(x'-z)^2}\mathcal C(z)\bigg],
\label{eq:Fc}
\end{align}
where $\mathcal C(z)$ is $\mathcal C(y')$ with $U^{e'c'}(y')\to U^{e'c'}(z)$, and similarly for $a$, $b$. $F_n^{\Delta}(p^+)$ falls faster than $\frac{1}{p^+}$:
\begin{itemize}
\item $G_n(y')-G_n(z)$ vanishes at $y'=z$, so the product with $\frac{(y'-z)^m}{(y'-z)^2}$ is bounded there, and its Fourier transform at $q=tk$ falls like $\frac{1}{q^2}\propto\frac{k^{+2}}{(p^++k^+)^2}$;
\item at $y'=x'$ the factor $\frac{(x'-y')^{k'}}{(x'-y')^2}$ gives a $\frac{1}{q}$ term, but with the phase $e^{-iq\cdot(x'-z)}$, which oscillates in $p^+$. Its $p^+$ integral converges.
\end{itemize}
So $\int^{\vee}dp^+\,\frac{F_n^\Delta}{k^+}$, $\int^\vee dp^+\frac{F_n(p^+)}{p^++k^+}$ and $\int^\vee dp^+\frac{F_n(p^+)}{p^+}$ are all finite for $\vee\to\infty$. The only $p^+$ integrand of the row that falls like $\frac{1}{p^+}$ is the $G(z)$ part of the $-\frac{1}{k^+}$ term:
\begin{equation}
-\frac{C}{k^+}\int_0^{\vee-k^+}\frac{dp^+}{2\pi}\,\frac{k^+}{p^++k^+}\,F_c^z=-\frac{C}{2\pi}\,F_c^z\,\log\frac{\vee}{k^+}\qquad\text{exactly}.
\label{eq:iii}
\end{equation}

\paragraph{Step 5: (iii), the term (B), written out.} Eq.~\eqref{eq:iii} with Eq.~\eqref{eq:Fc}:
\begin{align}
&-\frac{1}{(2\pi)^{3}}\, \frac{ig^{4}\,f^{c'ca}}{8\pi^{5}}\, \frac{1}{k^{+}}\,\log{\left(\frac{\vee}{k^{+}}\right)}\int_{w,w',z,y,x,x'} e^{-ik\cdot (w'-w)}\,\frac{(w'-z)\cdot (y-w)\,(x'-z)\cdot (x-z)}{(w'-z)^{2}\,(y-w)^{2}\,(x'-z)^{2}\,(x-z)^{2}}\notag \\
& \left[U^{e'c'}(z)\,\rho^{e'}(x')-U^{c'e'}(x')\,\rho^{e'}(x')\right]\left[U^{ce}(x)\,U^{ad}(y)\,\rho^{e}(x)\rho^{d}(y)-U^{ad}(w)\,U^{ce}(x)\,\rho^{e}(x)\rho^{d}(y)\right]\tag{(B)$_{\rm V}$}
\end{align}

This is the blue (iii) line of Group III Row V in \texttt{Checks\_JIMWLK\_BFKL.pdf}.

\paragraph{Step 6: the row.} Steps 3--5 give, exactly,
\begin{equation}
\frac{d^3N_{\rm V}}{d^3k}=\underbrace{-\frac{C}{2\pi}F_b(0)\log\frac{k^+}{\Lambda}}_{\text{JIMWLK}}\;+\;\underbrace{\text{(i)}+\text{(ii)}}_{=0}\;+\;\underbrace{\text{(B)}_{\rm V}}_{\text{BFKL}}\;+\;R_{\rm V},
\end{equation}
where $R_{\rm V}$ is line 1 with $\left[\frac{1}{p^+}\right]_+$ replaced as in Eq.~\eqref{eq:I2}, and with $G(y')\to G(y')-G(z)$ in the $-\frac{1}{k^+}$ term, plus line 2 with $\log(1-\frac{k^+}{\vee})$:
\begin{align}
R_{\rm V}&=  \frac{1}{(2\pi)^{3}}\, \frac{ig^{4}\,f^{c'ca}}{4\pi^{4}}\, \frac{1}{k^{+}}\,\int_{0}^{\vee-k^{+}}\, \frac{dp^{+}}{2\pi}\, \int_{w,y',z,y,x,x'}\, e^{-ik\cdot (y'-w)}\, e^{-ik\cdot \frac{p^{+}}{k^{+}}\,(y'-z)}\,\frac{(y'-z)^{m}\,(x'-y')^{k'}\,(y-w)^{i}\,(x-z)^{j}}{(y'-z)^{2}\,(x'-y')^{2}\,(y-w)^{2}\,(x-z)^{2}}\notag \\ & \left[\frac{\delta_{k'm}\,\delta_{ij}}{(p^{+}+k^{+})}-\old{\left[\frac{1}{p^{+}}\right]_{+}}\new{\left(\left[\frac{k^{+}}{p^{+}\,(p^{+}+k^{+})}\right]_{+}+\frac{1}{(p^{+}+k^{+})}\right)}\,\delta_{jm}\,\delta_{ik'}\old{-\frac{1}{k^{+}}\, \delta_{im}\, \delta_{jk'}}\right]\, \left[U^{e'c'}(y')\,\rho^{e'}(x')-U^{c'e'}(x')\,\rho^{e'}(x')\right]\times \notag \\ &  \left[U^{ce}(x)\,U^{ad}(y)\,\rho^{e}(x)\rho^{d}(y)-U^{ad}(w)\,U^{ce}(x)\,\rho^{e}(x)\rho^{d}(y)\right] \notag \\
& \nc  +\,\frac{1}{(2\pi)^{3}}\, \frac{ig^{4}\,f^{c'ca}}{4\pi^{4}}\, \frac{1}{k^{+}}\,\int_{0}^{\vee-k^{+}}\, \frac{dp^{+}}{2\pi}\, \int_{w,y',z,y,x,x'}\, e^{-ik\cdot (y'-w)}\, e^{-ik\cdot \frac{p^{+}}{k^{+}}\,(y'-z)}\,\frac{(y'-z)^{m}\,(y-w)^{i}\,(x-z)^{j}}{(y'-z)^{2}\,(y-w)^{2}\,(x-z)^{2}}\notag \\ & \nc  \left[-\frac{1}{k^{+}}\, \delta_{im}\, \delta_{jk'}\right]\, \Bigg[\frac{(x'-y')^{k'}}{(x'-y')^{2}}\,\left[U^{e'c'}(y')\,\rho^{e'}(x')-U^{c'e'}(x')\,\rho^{e'}(x')\right]\notag \\ & \nc  \qquad-\frac{(x'-z)^{k'}}{(x'-z)^{2}}\,\left[U^{e'c'}(z)\,\rho^{e'}(x')-U^{c'e'}(x')\,\rho^{e'}(x')\right]\Bigg]\times \notag \\ & \nc   \left[U^{ce}(x)\,U^{ad}(y)\,\rho^{e}(x)\rho^{d}(y)-U^{ad}(w)\,U^{ce}(x)\,\rho^{e}(x)\rho^{d}(y)\right] \notag \\
& -\,\frac{1}{(2\pi)^{3}}\, \frac{ig^{4}\,f^{c'ca}}{8\pi^{5}}\, \frac{1}{k^{+}}\,\log{\left(1-\frac{k^{+}}{\vee}\right)} \int_{w,y',z,y,x,x'}\, e^{-ik\cdot (y'-w)}\,\frac{(y'-z)\cdot (x-z)\,(x'-y')\cdot (y-w)}{(y'-z)^{2}\,(x'-y')^{2}\,(y-w)^{2}\,(x-z)^{2}}\notag \\ & \left[U^{e'c'}(y')\,\rho^{e'}(x')-U^{c'e'}(x')\,\rho^{e'}(x')\right]\times  \left[U^{ce}(x)\,U^{ad}(y)\,\rho^{e}(x)\rho^{d}(y)-U^{ad}(w)\,U^{ce}(x)\,\rho^{e}(x)\rho^{d}(y)\right]
\end{align}

Every $p^+$ integrand in $R_{\rm V}$ falls at least like $\frac{1}{p^{+2}}$ or oscillates, so $R_{\rm V}$ is finite for $\vee\to\infty$. ($\left[\frac{k^+}{p^+(p^++k^+)}\right]_+F\equiv\frac{k^+[F(p^+)-F(0)]}{p^+(p^++k^+)}$.) This is the uniqueness argument: had we put a different amount of $\log\frac{\vee}{k^+}$ into the BFKL list, $R$ would grow like $\log\vee$.

\paragraph{Numerical check} (\texttt{bfkl\_rows/exact\_check.py}; line 1 of your new file is identical to the one checked there). Two configurations, $\vee=20k^+$, the $p^+$ integral done numerically, pointwise in the transverse points:
\begin{itemize}
\item line 1 $=$ (rewritten line 1) $+$ ($G(y')-G(z)$ term) $+$ (ii) $+$ ($G(z)$ term), Rows III(1), III(2), V, VI: $\le3.7\times10^{-15}$;
\item the $G(z)$ term after the Fourier identity \eqref{eq:I3} and the $p^+$ integral $=$ (B): $\le3.4\times10^{-15}$;
\item the Fourier transforms themselves: $q\int_0^\infty J_1(qr)\,e^{-ar}dr$ at $a=10^{-2},3\times10^{-3},10^{-3}$ (must tend to 1): $q=0.7$: 0.9857, 0.9957, 0.9986; $q=1.9$: 0.9947, 0.9984, 0.9995. $2\pi\int_0^\infty rJ_0(qr)\log r^2\,e^{-ar^2}dr\big/(-\frac{4\pi}{q^2})$ at $a=4,2,1\times10^{-3}$, extrapolated to $a=0$ (must be 1): $q=0.7$: 1.0350, 1.0169, 1.0083 $\to$ 0.99971; $q=1.9$: 1.0045, 1.0022, 1.0011 $\to$ 1.00000.
\end{itemize}

\section{Group III Rows VI and III (first part)}
Your Eqs.~(1.15)--(1.17) (Row VI) and (1.5)--(1.7) (Row III part 1) have the same three terms in the bracket. The steps are those of Sec.~\ref{sec:V} word for word. Only the colour factor $\mathcal C(y')$ changes; it is $\big[U^{e'c'}(y')U^{d'b}(z)\rho^{e'}(x')-U^{bd'}(z)U^{c'e'}(x')\rho^{e'}(x')\big]\times[\dots]$ (and $f^{c'd'a}$). The results:
\begin{align}
&-\frac{1}{(2\pi)^{3}}\, \frac{ig^{4}\,f^{c'd'a}}{8\pi^{5}}\, \frac{1}{k^{+}}\,\log{\left(\frac{\vee}{k^{+}}\right)}\int_{w,w',z,y,x,x'} e^{-ik\cdot (w'-w)}\,\frac{(w'-z)\cdot (y-w)\,(x'-z)\cdot (x-z)}{(w'-z)^{2}\,(y-w)^{2}\,(x'-z)^{2}\,(x-z)^{2}}\notag \\
& \left[U^{e'c'}(z)\,U^{d'b}(z)\,\rho^{e'}(x')-U^{bd'}(z)\,U^{c'e'}(x')\,\rho^{e'}(x')\right]\notag \\
& \times\left[U^{bc}(z)\,U^{ad}(y)\,U^{ce}(x)\,\rho^{d}(y)\,\rho^{e}(x)-U^{ad}(y)\,\rho^{d}(y)\,\rho^{b}(x)\right]\tag{(B)$_{\rm VI}$}
\end{align}

\begin{align}
&-\frac{1}{(2\pi)^{3}}\, \frac{ig^{4}\,f^{c'd'a}}{16\pi^{5}}\, \frac{1}{k^{+}}\,\log{\left(\frac{\vee}{k^{+}}\right)}\int_{w,x',w',z,y,x} e^{-ik\cdot (w'-w)}\,\frac{(w'-z)\cdot (y-w)\,(x'-z)\cdot (x-z)}{(w'-z)^{2}\,(y-w)^{2}\,(x'-z)^{2}\,(x-z)^{2}}\notag \\
& \left[U^{e'c'}(z)\,U^{d'b}(z)\,\rho^{e'}(x')-U^{bd'}(z)\,U^{c'e'}(x')\,\rho^{e'}(x')\right]\notag \\
& \times\left[U^{ac}(w)\, \{\rho^{c}(y),\rho^{b}(x)\}-U^{bc}(z)\,U^{ad}(y)\,U^{ce}(x)\,\{\rho^{d}(y),\rho^{e}(x)\}\right]\tag{(B)$_{\rm III(1)}$}
\end{align}

\noindent (i) and (ii) cancel as in Eq.~\eqref{eq:i_ii}.

\section{Group III Row III, second part}
Your Eq.~(1.10), line 1:
\begin{align}
&\frac{1}{(2\pi)^{3}}\,\frac{g^{4}}{4\pi^{4}}\,\int_{0}^{\vee-k^{+}}\, \frac{dp^{+}}{2\pi}\,\int_{w,w',y',z,y,x}\, e^{-ik\cdot (y'-w)}\,e^{-ik\cdot \frac{p^{+}}{k^{+}}\left(y'-z\right)}\, \frac{(y'-z)^{m}\,(x'-y')^{k'}\,}{(y'-z)^{2}\,(x'-y')^{2}\,}\, \frac{1}{\left[p^{+}\,(x-z)^{2}+k^{+}\,(x-w)^{2}\right]}\notag \\ & \Bigg[\frac{\delta_{k'm}\,d_{\perp}\,p^{+}}{2\,(z-w)^{2}\,(p^{+}+k^{+})^{2}}\, \left[(x-z)^{2}-(x-w)^{2}\right]+\frac{p^{+}\,\delta_{k'm}}{k^{+}\,(p^{+}+k^{+})}\,\bigg[\frac{(x-z)\cdot (z-w)}{(z-w)^{2}} -\frac{(x-z)\cdot (x-w)}{2\,(x-w)^{2}}\notag \\ & -\frac{\left[(x-z)^{2}-(x-w)^{2}\right]}{2\,(z-w)^{2}}\bigg] +\frac{\delta_{k'm}}{(p^{+}+k^{+})}\,\bigg[\frac{(x-w)\cdot (z-w)}{(z-w)^{2}}+\frac{(x-z)\cdot (x-w)}{2\,(x-z)^{2}} -\frac{\left[(x-z)^{2}-(x-w)^{2}\right]}{2\,(z-w)^{2}}\bigg]\notag \\ & -\left[\frac{1}{p^{+}}\right]_{+}\,\bigg[\frac{(x-w)^{k'}\,(z-w)^{m}}{(z-w)^{2}}+\frac{(x-z)^{m}\,(x-w)^{k'}}{2\,(x-z)^{2}}\bigg]-\frac{p^{+}}{k^{+\,2}}\,\bigg[\frac{(x-z)^{k'}\,(z-w)^{m}}{(z-w)^{2}} -\,\frac{(x-z)^{k'}\,(x-w)^{m}}{2\,(x-w)^{2}}\bigg]\notag \\ & -\frac{1}{k^{+}}\,\bigg[\frac{(x-z)^{m}\,(z-w)^{k'}}{(z-w)^{2}}-\frac{(x-z)^{m}\,(x-w)^{k'}}{2\,(x-w)^{2}}+\frac{(x-w)^{m}\,(z-w)^{k'}}{(z-w)^{2}} +\frac{(x-z)^{k'}\,(x-w)^{m}}{2\,(x-z)^{2}}\bigg]\Bigg]\times\,\notag \\ & \left[f^{c'd'a}\,U^{e'c'}(y')\,U^{d'b}(z)\,\rho^{e'}(x')-f^{c'd'a}\, 
    U^{bd'}(z)\,U^{c'e'}(x')\,\rho^{e'}(x')\right]\times \left[f^{cbd}\,U^{ac}(w)\,\rho^{d}(x) -f^{acd}\,U^{bc}(z)\,U^{de}(x)\,\rho^{e}(x)\right]\tag{your (1.10), line 1}
\end{align}

The bracket has six terms. After the $y'$ integral, Eq.~\eqref{eq:Fsplit}, each $G(z)$ part gets $\frac{k^+}{p^++k^+}$; the common factor is $\frac{1}{p^+(x-z)^2+k^+(x-w)^2}\sim\frac{1}{p^+}$. The large-$p^+$ power of each term, exactly:
\begin{center}\renewcommand{\arraystretch}{1.5}
\begin{tabular}{lll}\toprule
term & coefficient $\times\frac{1}{p^+A+k^+B}\times\frac{k^+}{p^++k^+}$ & $\int^{\vee}dp^+$\\\midrule
$\frac{\delta_{k'm}d_\perp p^+}{2(z-w)^2(p^++k^+)^2}[\dots]$ & $\sim p^{+-3}$ & finite\\
$\frac{p^+\delta_{k'm}}{k^+(p^++k^+)}[\dots]$ & $\sim p^{+-2}$ & finite\\
$\frac{\delta_{k'm}}{p^++k^+}[\dots]$ & $\sim p^{+-3}$ & finite\\
$-\left[\frac{1}{p^+}\right]_+[\dots]$ & $F(p^+)$ part $\sim p^{+-3}$; $-F(0)$ part: (ii), as in Eq.~\eqref{eq:I2} & (ii) $=-$(i)\\
$-\frac{p^+}{k^{+2}}[\dots]$ & $\sim p^{+-1}$ & $\log\frac{\vee}{k^+}$: (iii)\\
$-\frac{1}{k^+}[\dots]$ & $\sim p^{+-2}$ & finite\\\bottomrule
\end{tabular}
\end{center}
with $A=(x-z)^2$, $B=(x-w)^2$. For the $-\frac{p^+}{k^{+2}}$ term, the partial fractions (I5) are exact:
\begin{equation}
-\frac{p^+}{k^{+2}}\,\frac{k^+}{p^++k^+}\,\frac{1}{p^+A+k^+B}=-\frac{1}{k^+A}\,\frac{1}{p^++k^+}+\frac{B}{A\,(p^++k^+)(p^+A+k^+B)}.
\label{eq:I5}
\end{equation}
The first term gives $\int_0^{\vee-k^+}\frac{dp^+}{p^++k^+}=\log\frac{\vee}{k^+}$, i.e.\ (B); the second falls like $\frac{1}{p^{+2}}$ and stays in the remaining terms. Written out:
\begin{align}
&-\frac{1}{(2\pi)^{3}}\,\frac{g^{4}}{8\pi^{5}}\,\frac{1}{k^{+}}\,\log{\left(\frac{\vee}{k^{+}}\right)}\int_{w,x',w',z,x} e^{-ik\cdot (w'-w)}\,\frac{(x'-z)\cdot (x-z)}{(x'-z)^{2}\,(x-z)^{2}}\notag \\
& \times\bigg[\frac{(w'-z)\cdot (z-w)}{(w'-z)^{2}\,(z-w)^{2}}-\frac{(w'-z)\cdot (x-w)}{2\,(w'-z)^{2}\,(x-w)^{2}}\bigg]\notag \\
& \left[f^{c'd'a}\,U^{e'c'}(z)\,U^{d'b}(z)\,\rho^{e'}(x')-f^{c'd'a}\, U^{bd'}(z)\,U^{c'e'}(x')\,\rho^{e'}(x')\right]\notag \\
& \times\left[f^{cbd}\,U^{ac}(w)\,\rho^{d}(x) -f^{acd}\,U^{bc}(z)\,U^{de}(x)\,\rho^{e}(x)\right]\tag{(B)$_{\rm III(2)}$}
\end{align}

In the remaining terms (\texttt{Remaining\_RC\_DGLAP\_finite.pdf}, Group III Row III part 2) the $-\frac{p^+}{k^{+2}}$ term becomes its $G(y')-G(z)$ part, and the second term of Eq.~\eqref{eq:I5} is added:
\begin{align}
&\nc +\,\frac{1}{(2\pi)^{3}}\,\frac{g^{4}}{4\pi^{4}}\,\int_{0}^{\vee-k^{+}}\, \frac{dp^{+}}{2\pi}\,\int_{w,x',w',z,x}\, e^{-ik\cdot (w'-w)}\, \frac{(w'-z)^{m}\,(x'-z)^{k'}}{(w'-z)^{2}\,(x'-z)^{2}}\, \notag \\ & \nc  \Bigg[ \frac{(x-w)^{2}}{(x-z)^{2}\,(p^{+}+k^{+})\,\left[p^{+}\,(x-z)^{2}+k^{+}\,(x-w)^{2}\right]}\,\bigg[\frac{(x-z)^{k'}\,(z-w)^{m}}{(z-w)^{2}} -\,\frac{(x-z)^{k'}\,(x-w)^{m}}{2\,(x-w)^{2}}\bigg] \Bigg]\times\,\notag \\ & \nc  \left[f^{c'd'a}\,U^{e'c'}(z)\,U^{d'b}(z)\,\rho^{e'}(x')-f^{c'd'a}\, 
    U^{bd'}(z)\,U^{c'e'}(x')\,\rho^{e'}(x')\right]\times \left[f^{cbd}\,U^{ac}(w)\,\rho^{d}(x) -f^{acd}\,U^{bc}(z)\,U^{de}(x)\,\rho^{e}(x)\right]
\end{align}


\section{Group II Row IV}\label{sec:G2IV}
\paragraph{Step 1: your equations.} With $\xi=\frac{k^+}{p^++k^+}$, your Eq.~(1.19), and after the plus prescription, your Eq.~(1.21):
\begin{align}
    \frac{d^{3}N_{\rm NLO}^{\rm}}{d^{3}k}  & = \frac{g^{2}}{4\pi^{2}}\, \frac{N_{c}}{(2\pi)^{3}}\,\frac{g^{2}}{2\pi^{2}}\, \frac{1}{k^{+}}\,\int_{x,y,x',y'}\,\frac{(x-y)\cdot (x'-y')}{(x-y)^{2}\,(x'-y')^{2}}\,\left[-\frac{d_{\perp}}{\epsilon}-d_{\perp}\,\gamma_{E}-\frac{d_{\perp}}{2}\,\log{\left(\frac{(y'-y)^{2}\,\mu^{2}_{\overline{MS}}}{4}\right)}-1\right]\,\notag \\ & \Bigg[\int_{\frac{k^{+}}{\vee}}^{1}\,\frac{d\xi}{\xi^{2}}\,e^{-ik\cdot (y'-y)\,\frac{1}{\xi}}\,\left[\xi\,(1-\xi)+\frac{(1-\xi)}{\xi}+\frac{\xi}{(1-\xi)_{+}}\right]\Bigg]\times \left[U^{e'c}(y') - U^{e'c}(x')\right] \times\, \left[U^{ec}(y) - U^{ce}(x)\right]\, \rho^{e'}(x') \rho^{e}(x)\notag \\ &  +\, \frac{g^{2}}{4\pi^{2}}\, \frac{1}{(2\pi)^{3}}\,\frac{g^{2}}{2\pi^{2}}\, \frac{N_{c}}{k^{+}}\,\left[\log{\left(\frac{k^{+}}{\Lambda}\right)}+\log{\left(\frac{\vee}{k^{+}}-1\right)}\right]\,\int_{x,y,x',y'}\,e^{-ik\cdot (y'-y)}\,\frac{(x-y)\cdot (x'-y')}{(x-y)^{2}\,(x'-y')^{2}}\,\bigg[-\frac{d_{\perp}}{\epsilon}-d_{\perp}\,\gamma_{E}\notag \\ & -\frac{d_{\perp}}{2}\,\log{\left(\frac{(y'-y)^{2}\,\mu^{2}_{\overline{MS}}}{4}\right)}-1\bigg]\,\left[U^{e'c}(y') - U^{e'c}(x')\right] \times\, \left[U^{ec}(y) - U^{ce}(x)\right]\, \rho^{e'}(x') \rho^{e}(x)\notag \\ & +\,\frac{g^{2}}{4\pi^{2}}\, \frac{1}{(2\pi)^{3}}\,\frac{g^{2}}{2\pi^{2}}\, \frac{N_{c}}{k^{+}}\,\int_{\frac{k^{+}}{\vee}}^{\frac{k^{+}}{k^{+}+\Lambda}}\,d\xi\, \frac{(1-\xi)}{\xi}\, \int_{x,y,x',y'}\,e^{-ik\cdot (y'-y)\frac{1}{\xi}}\,\Bigg[\frac{(x-y)\cdot (x'-y')}{(x-y)^{2}\, (x'-y')^{2}} -\frac{(x-y)\cdot (y'-y)\, (x'-y')\cdot (y'-y)}{(x-y)^{2}\,(x'-y')^{2}\,(y'-y)^{2}}\,d_{\perp}\Bigg]\times\notag \\ &  \left[U^{e'c}(y') - U^{e'c}(x')\right] \times\, \left[U^{ec}(y) - U^{ce}(x)\right]\, \rho^{e'}(x') \rho^{e}(x)\tag{your (1.21)}
\end{align}

Call $g(\xi)$ the $\xi$-dependent part of the integrand of line 1 at fixed transverse points: $g(\xi)\propto e^{-ik\cdot(y'-y)\frac1\xi}$. Line 2 is $g(1)$ times $\int_{\frac{k^+}{\vee}}^{\frac{k^+}{k^++\Lambda}}\frac{d\xi}{\xi(1-\xi)}=\log\frac{k^+}{\Lambda}+\log(\frac{\vee}{k^+}-1)$: your plus prescription, Eq.~\eqref{eq:plus}, in the $\xi$ form of Eq.~\eqref{eq:step3}.

\paragraph{Step 2: (i) and (ii).} $\frac{1}{\xi(1-\xi)}=\frac{1}{\xi}+\frac{1}{1-\xi}$ and $\int_{\frac{k^+}{\vee}}^1\frac{d\xi}{\xi}=\log\frac{\vee}{k^+}$ give, exactly,
\begin{equation}
\int_{\frac{k^+}{\vee}}^1\frac{d\xi}{\xi^2}\frac{\xi}{(1-\xi)_+}g=\int_{\frac{k^+}{\vee}}^1\frac{d\xi}{\xi^2}\Big(\xi+\frac{\xi^2}{(1-\xi)_+}\Big)g\;-\;g(1)\log\frac{\vee}{k^+},\qquad \int\frac{d\xi}{\xi^2}\frac{\xi^2}{(1-\xi)_+}g\equiv\int\frac{d\xi}{1-\xi}\big[g(\xi)-g(1)\big].
\label{eq:I2xi}
\end{equation}
The last term, times the prefactor of line 1, is (ii). It is your \texttt{BFKL.tex} line 1 for this row; it is exactly minus (i), the $\log\frac{\vee}{k^+}$ of line 2, poles included ($5.5\times10^{-16}$).

\paragraph{Step 3: (iii), $H_4$.} Which $\xi$ integrand of line 1 falls like $\frac{1}{\xi}$ for $\xi\to\frac{k^+}{\vee}\to0$? The $y'$ integral is a Fourier transform at $q=\frac{k}{\xi}$, large for small $\xi$. Write $H(y,y')=\frac{(x-y)\cdot(x'-y')}{(x-y)^2(x'-y')^2}[U^{e'c}(y')-U^{e'c}(x')][U^{ec}(y)-U^{ce}(x)]\rho^{e'}(x')\rho^e(x)=H(y,y)+[H(y,y')-H(y,y)]$.
\begin{itemize}
\item The constants of the pole bracket times $H(y,y)$: $\int_{y'}e^{-ik\cdot(y'-y)\frac1\xi}=(2\pi)^2\xi^2\delta^{(2)}(k)=0$ for $k\neq0$.
\item $-\frac{d_\perp}{2}\log\frac{(y'-y)^2\mu^2}{4}$ times $H(y,y)$: $\int_{y'}e^{-ik\cdot(y'-y)\frac1\xi}\log\frac{(y'-y)^2\mu^2}{4}=-\frac{4\pi\xi^2}{k^2}$ ($\mu$ drops out for $k\neq0$). With the measure $\frac{d\xi}{\xi^2}$, the three terms of the bracket give $\int d\xi\,\xi(1-\xi)$, $\int d\xi\,\xi$ (the $\xi$ of Eq.~\eqref{eq:I2xi}), $\int d\xi\frac{\xi}{1-\xi}\cdot$(subtracted): finite; and
\begin{equation}
\int_{\frac{k^+}{\vee}}^1d\xi\,\frac{1-\xi}{\xi}=\log\frac{\vee}{k^+}-1+\frac{k^+}{\vee}.
\end{equation}
\item $H(y,y')-H(y,y)$ vanishes at $y'=y$: its transform falls faster, finite.
\end{itemize}
So the only $\log\frac{\vee}{k^+}$ of line 1 besides (ii) is the $\frac{1-\xi}{\xi}$ term with $\log(y'-y)^2$ at $y'=y$. With $\frac{4\pi}{k^2}=\frac{1}{\pi}\int_{w,w'}e^{-ik\cdot(w'-w)}\Kc(w-z)\cdot\Kc(w'-z)$ and $y\to z$, its $\log\frac{\vee}{k^+}$ term is, exactly (identity (I4)),
\begin{align}
&\frac{1}{(2\pi)^{3}}\,\frac{g^{4}\,N_{c}}{8\pi^{5}}\,\frac{1}{k^{+}}\,\log{\left(\frac{\vee}{k^{+}}\right)}\,\int_{x,x',z,w,w'}\,e^{-ik\cdot (w'-w)}\,\frac{(w-z)\cdot(w'-z)}{(w-z)^{2}\,(w'-z)^{2}}\,\frac{(x-z)\cdot(x'-z)}{(x-z)^{2}\,(x'-z)^{2}}\notag \\ & \left[U^{e'c}(z) - U^{e'c}(x')\right] \times\, \left[U^{ec}(z) - U^{ce}(x)\right]\, \rho^{e'}(x') \rho^{e}(x)\tag{$\log\frac{\vee}{k^+}H_4$}
\end{align}

your \texttt{BFKL.tex} line 2 of this row. Its rest, $\big(-1+\frac{k^+}{\vee}\big)H_4\big|_{\log\to1}$, stays in the remaining terms. ($d_\perp=2$ here: no pole multiplies it.)

\paragraph{Step 4: the row.} Group II Row IV $=\log\frac{k^+}{\Lambda}$ (your line 2, JIMWLK) $+$ (i) $+$ (ii) $+\log\frac{\vee}{k^+}H_4$ $+$ a rest that is finite for $\vee\to\infty$; (i) $+$ (ii) $=0$.

\section*{Part B. The coefficient in front of $\log\mu^2$: running coupling and DGLAP}
\section{The five sources of \texorpdfstring{$\mu$}{mu}}\label{sec:mu}
As your advisor says, $\mu$ enters through A3 and the shockwave through the loop, each with its c.c., and through the real term: five sources. Each comes in the combination $\hat E_r=\frac{2}{\epsilon}+\log\frac{r^2\mu^2}{c_0^2}$ ($c_0=2e^{-\gamma_E}$), so the coefficient of $\log\mu^2$ is the coefficient of $\frac{2}{\epsilon}$. In units of
\begin{equation*}
{\rm LO}\equiv\frac{1}{(2\pi)^{3}}\frac{g^{2}}{2\pi^{2}}\frac{1}{k^{+}}\int_{x,y,x',y'}e^{-ik\cdot(y'-y)}\frac{(x-y)\cdot(x'-y')}{(x-y)^{2}(x'-y')^{2}}\left[U^{e'c}(y')-U^{e'c}(x')\right]\left[U^{ec}(y)-U^{ce}(x)\right]\rho^{e'}(x')\rho^{e}(x)
\end{equation*}
($\alpha_s=\frac{g^2}{4\pi}$; $L\equiv\log\frac{k^+}{\Lambda}$, $l_1\equiv\log(1-\frac{k^+}{\vee})$, $X=x-y$, $\Delta=y'-y$):
\paragraph{A3, your Group I Eq.~(1.2).} Prefactor $\frac{1}{(2\pi)^3}\frac{g^4N_c}{32\pi^4}\frac{1}{k^+}=\frac{\alpha_sN_c}{4\pi}\times$(LO prefactor). With $w'\to y'$, $w\to y$ the colour factor is $-$(that of LO) and the bracket is $-\hat E_X$:
\begin{equation}
\text{A3}+\text{c.c.}=\frac{\alpha_sN_c}{4\pi}\,{\rm LO}\Big[(\hat E_X+\hat E_{X'})\Big(\frac{11}{3}-2L+l_1\Big)\Big]+\dots\quad\Rightarrow\quad\log\mu^2:\ \frac{\alpha_s}{\pi}\,\frac{N_c}{2}\Big(\frac{11}{3}-2L+l_1\Big){\rm LO}.
\label{eq:muA3}
\end{equation}
\paragraph{Shockwave through the loop, your Group I Eq.~(1.32).} Prefactor $\frac{1}{(2\pi)^3}\frac{g^4N_c}{16\pi^4}\frac{1}{k^+}=\frac{\alpha_sN_c}{2\pi}\times$(LO prefactor). With $y\to x$, $z\to w$ (dummy variables) the phase, kernel and colour factor are those of A3 with the opposite sign of the colour factor, and the bracket is $+\hat E$:
\begin{equation}
\text{SW}+\text{c.c.}=\frac{\alpha_sN_c}{2\pi}\,{\rm LO}\Big[(\hat E_X+\hat E_{X'})\Big(-\frac{11}{6}+2L+2\log(1-\tfrac{\Lambda}{k^+})\Big)\Big]\quad\Rightarrow\quad\log\mu^2:\ \frac{\alpha_s}{\pi}N_c\Big(-\frac{11}{6}+2L\Big){\rm LO}.
\label{eq:muSW}
\end{equation}
\paragraph{Real, your Group II Eq.~(1.21).} Prefactor $\frac{g^2}{4\pi^2}N_c\times$(LO prefactor)$=\frac{\alpha_s}{\pi}N_c$. The bracket $\big[-\frac{d_\perp}{\epsilon}-d_\perp\gamma_E-\frac{d_\perp}{2}\log\frac{(y'-y)^2\mu^2}{4}-1\big]=-\hat E_\Delta+O(\epsilon)$ ($d_\perp=2-\epsilon$). Line 1 with your $\frac{\xi}{(1-\xi)_+}$ is, by Eq.~\eqref{eq:I2xi} and the $[0,1]$ plus distribution restricted to $[\frac{k^+}{\vee},1]$ ($\int_a^1\frac{d\xi\,h}{(1-\xi)_+}=\int_a^1\frac{h(\xi)-h(1)}{1-\xi}d\xi+h(1)\log(1-a)$, the $\overline{MS}$ one), $P^{(r)}_{gg,+}\otimes{\rm LO}-2N_c\big(l_1+\log\frac{\vee}{k^+}\big){\rm LO}$:
\begin{align}
\text{line 1}&\Rightarrow\ \log\mu^2:\ -\frac{\alpha_s}{2\pi}\Big[\int_{k^+/\vee}^1\frac{d\xi}{\xi^2}P^{(r)}_{gg,+}(\xi)\,{\rm LO}_\xi-2N_c\Big(l_1+\log\frac{\vee}{k^+}\Big){\rm LO}\Big],\label{eq:muR1}\\
\text{line 2}&\Rightarrow\ \log\mu^2:\ -\frac{\alpha_s}{2\pi}\,2N_c\Big(L+\log\frac{\vee}{k^+}+l_1\Big){\rm LO}.\label{eq:muR2}
\end{align}
$P^{(r)}_{gg,+}=2N_c\big[\frac{\xi}{(1-\xi)_+}+\frac{1-\xi}{\xi}+\xi(1-\xi)\big]$ is the real part of $P_{gg}$; it has no $\delta(1-\xi)$.
\paragraph{The coupling of LO.} LO is $\propto g^2$. In $\overline{MS}$, $g^2=g^2(\mu_R)\big[1-\frac{\alpha_s\beta_0}{4\pi}\hat E\big]$, $\beta_0=\frac{11}{3}N_c$:
\begin{equation}
\text{CT}=-\frac{\alpha_s\beta_0}{4\pi}\,\hat E\;{\rm LO}\quad\Rightarrow\quad\log\mu^2:\ -\frac{\alpha_s}{\pi}\,\frac{\beta_0}{4}\,{\rm LO}.
\label{eq:ct}
\end{equation}

\paragraph{Sum them up.} Coefficient of $\log\mu^2$ in units of $\frac{\alpha_s}{\pi}{\rm LO}$ (\texttt{mu\_check.py}; $P^*\equiv\int\frac{d\xi}{\xi^2}P^{(r)}_{gg,+}{\rm LO}_\xi/{\rm LO}$):
{\small
\begin{longtable}{@{}p{0.36\textwidth}cccccc@{}}\toprule
source & $\beta_0$ ($\frac{11}{3}$) & $\times L$ & $\times\log\frac{\vee}{k^+}$ & $\times l_1$ & $\times P^*$\\\midrule
A3, Group I Row II part 1 $+$ c.c., Eq.~\eqref{eq:muA3} & $11N_c/6$ & $-N_c$ & $0$ & $N_c/2$ & $0$\\
shockwave through the loop, Group I Row V 2b UV $+$ c.c., Eq.~\eqref{eq:muSW} & $-11N_c/6$ & $2N_c$ & $0$ & $0$ & $0$\\
real, Group II Row IV line 1, Eq.~\eqref{eq:muR1} & $0$ & $0$ & $N_c$ & $N_c$ & $-1/2$\\
real, Group II Row IV line 2, Eq.~\eqref{eq:muR2} & $0$ & $-N_c$ & $-N_c$ & $-N_c$ & $0$\\
coupling of LO, Eq.~\eqref{eq:ct} & $-11N_c/12$ & $0$ & $0$ & $0$ & $0$\\
\midrule
sum of the five sources & $0$ & $0$ & $0$ & $\frac{N_c}{2}$ & $-\frac12$\\
$+$ coupling of LO & $-\frac{11N_c}{12}=-\frac{\beta_0}{4}$ & $0$ & $0$ & $\frac{N_c}{2}$ & $-\frac12$\\\bottomrule
\end{longtable}}
\noindent So:
\begin{enumerate}
\item \textbf{$\beta_0$: the four sources of the single-gluon final state cancel exactly}: $\frac{N_c}{2}\cdot\frac{11}{3}-N_c\cdot\frac{11}{6}=0$. Not only $\log\mu^2$: the pole, $\gamma_E$ and $\log(x-w)^2$ cancel too, pointwise after the renaming (\texttt{dglap\_rc/checks.json}: $0$). This is what your advisor expected.
\item \textbf{$\times\log\frac{k^+}{\Lambda}$ (evolution): three sources}, A3 : shockwave : real $=-N_c:2N_c:-N_c$, sum $0$. They cancel only with the real term. Pointwise the poles are in the ratio $1:-2:1$ and cancel to $8.2\times10^{-16}$ (\texttt{ll\_poles.json}). These are the three sources of the evolution.
\item \textbf{$\times\log\frac{\vee}{k^+}$}: $+N_c$ (line 1, (ii)) and $-N_c$ (line 2, (i)) cancel: the BFKL terms have no $\mu$.
\item What is left: $-\frac{\alpha_s}{2\pi}\log\mu^2\big[\int P^{(r)}_{gg,+}\otimes{\rm LO}+\frac{\beta_0}{2}{\rm LO}\big]=-\frac{\alpha_s}{2\pi}\log\mu^2\,P_{gg}\otimes{\rm LO}$ (DGLAP, with $P_{gg}=P^{(r)}_{gg,+}+\frac{\beta_0}{2}\delta(1-\xi)$), and $+\frac{\alpha_s}{2\pi}N_c\,l_1\log\mu^2\,{\rm LO}$ (Q1, from A3 only; $l_1\to0$ for $\vee\to\infty$).
\end{enumerate}

\section{Why \texorpdfstring{$\beta_0$}{beta0} is not in the NLO terms, and whether Group II Row IV is running coupling}\label{sec:why}
\begin{itemize}
\item \textbf{The NLO diagrams carry no $\beta_0\log\mu^2$}, because of item 1: A3 ($+\frac{11}{3}$) and the shockwave through the loop ($2\times(-\frac{11}{6})$) cancel, and the real term has no $\beta_0$. So at this order nothing in your rows fixes a scale for the coupling; in particular no $\beta_0\log(x-w)^2$ survives.
\item \textbf{$\beta_0$ is still in the result: through the coupling of LO.} LO is $O(g^2)$, so $g^2$ must be renormalized; Eq.~\eqref{eq:ct} is the only $\beta_0$. Its pole is the $\frac{\beta_0}{2}\delta(1-\xi)$ part of $P_{gg}$, which the real term lacks. After $\overline{MS}$ renormalization of $g^2$ at $\mu_R$ and of $D_{h/g}$ at $\mu_F$, the $\beta_0$ term of the result is
\begin{equation}
\frac{\alpha_s\beta_0}{4\pi}\,\log\frac{\mu_R^2}{\mu_F^2}\;{\rm LO}\Big|_{\alpha_s(\mu_R)},
\end{equation}
and the result is $\mu_R$-independent: $\frac{d}{d\log\mu_R^2}{\rm LO}|_{\alpha_s(\mu_R)}=-\frac{\alpha_s\beta_0}{4\pi}{\rm LO}$. This is the running coupling: the coupling of LO runs, the NLO diagrams do not add to it.
\item \textbf{Group II Row IV is not running coupling.} Its $\log\mu^2$ (Eqs.~\eqref{eq:muR1}, \eqref{eq:muR2}) multiplies $\int\frac{d\xi}{\xi^2}P^{(r)}_{gg,+}(\xi){\rm LO}_\xi$, a convolution in $\xi$ with the \emph{real} splitting function. It has no $\delta(1-\xi)$ part and no $\beta_0$: the plus distribution integrates the $\frac{\xi}{1-\xi}$ singularity against $g(\xi)-g(1)$, and its $\delta(1-\xi)$ coefficient is zero in the $\overline{MS}$ convention, Eq.~\eqref{eq:muR1}. A running-coupling term would be $\beta_0\log\mu^2\times{\rm LO}$ (local in $\xi$). What Group II Row IV gives is the $\mu_F$ dependence of the fragmentation function (DGLAP), and $-2N_cL\log\mu^2$, which cancels against Group I (item 2).
\item With both: $-\frac{\alpha_s}{2\pi}\log\mu^2\,P_{gg}\otimes{\rm LO}$ is exactly what $\mu_F^2\frac{\partial}{\partial\mu_F^2}D_{h/g}=\frac{\alpha_s}{2\pi}P_{gg}\otimes D_{h/g}$ cancels. The coupling counterterm completes $P^{(r)}_{gg,+}$ to the full $P_{gg}$.
\end{itemize}
\end{document}
